% !TEX root = ../../../proposal.tex

\section{Diffusion Models}
\label{sec:diffusion_models}

Diffusion models \citep{pmlr-v37-sohl-dickstein15} are latent-variable generative models of the form
\(
p_\theta(x_0)
=
\int p_\theta(x_{0:T})\, \mathrm{d}x_{1:T},
\)
where \(x_{1:T}\) are latents of the same dimensionality as the observed data
\(x_0\). Following \citet{NEURIPS2020_4c5bcfec}, the model consists of a fixed
\emph{forward diffusion process} that gradually adds noise to data, and a learned
\emph{reverse process} that aims to invert this corruption.

\begin{figure}[H]
    \centering
    \includegraphics[width=0.65\textwidth]{chapters/background/media/DDPM.png}
    \caption{Directed graphical model structure of DDPMs \citep{NEURIPS2020_4c5bcfec}.}
    \label{fig:ddpm_graphical_model}
\end{figure}

The joint distribution of all variables is written as a Markov chain:
\begin{equation}
    p_\theta(x_{0:T})
    =
    p(x_T)
    \prod_{t=1}^{T} p_\theta(x_{t-1}\mid x_t),
    \qquad
    p(x_T)=\mathcal{N}(0,I),
    \label{eq:ddpm_joint}
\end{equation}
where \(p_\theta(x_{t-1}\mid x_t)\) is the reverse process. The forward diffusion
process is defined as a separate Markov chain:
\begin{equation}
    q(x_{1:T}\mid x_0)
    =
    \prod_{t=1}^T
    q(x_t\mid x_{t-1}),
    \qquad
    q(x_t\mid x_{t-1})
    =
    \mathcal{N}\!\left(
        \sqrt{1-\beta_t}\,x_{t-1},\;
        \beta_t I
    \right),
    \label{eq:ddpm_forward}
\end{equation}
where \(\{\beta_t\}_{t=1}^T\) is the variance schedule that \emph{can} be learned by reparameterization,
or fixed to predetermined constants (e.g., a linear schedule from $10^{-4}$ to $10^{-2}$) such that \(\beta_t\) typically increases slowly with \(t\).
This ensures that the forward chain remains
analytically tractable and preserves the same functional form across timesteps while gradually
driving \(x_T\) toward an isotropic Gaussian.

% \citet{NEURIPS2020_4c5bcfec} choose \(T\) large (e.g.\ \(T=1000\)) and define a monotonically increasing variance schedule \(\{\beta_t\}_{t=1}^T\), typically ranging from approximately \(10^{-4}\) to \(2\times10^{-2}\). 
% This yields a sequence of small Gaussian perturbations that progressively reduce the signal-to-noise ratio, such that \(x_T\) is approximately distributed as a standard Gaussian \(\mathcal{N}(0, I)\).
%%%%%%%%%%%%%%%%%%%%%%%%%%%%%%%%%%%%%%%%%%%%%%%%%%%%%%%%%%%%%%%%%%%%%%%%%%%%%%%
\subparagraph{Forward Diffusion Process.}
%%%%%%%%%%%%%%%%%%%%%%%%%%%%%%%%%%%%%%%%%%%%%%%%%%%%%%%%%%%%%%%%%%%%%%%%%%%%%%%

The forward diffusion process is a fixed Markov chain that gradually corrupts a clean sample \(x_0\) by adding Gaussian noise. Following Equation~\ref{eq:ddpm_forward}, each transition has the form
\begin{equation}
    q(x_t\mid x_{t-1})
    =
    \mathcal{N}\!\left(
        \sqrt{1-\beta_t}\,x_{t-1},\;
        \beta_t I
    \right).
    \label{eq:ddpm_forward_transition}
\end{equation}
Because these transitions are linear Gaussian, the marginal distribution at any timestep is available in closed form:
\begin{equation}
    q(x_t\mid x_0)
    =
    \mathcal{N}\!\left(
        \sqrt{\bar\alpha_t}\,x_0,\;
        (1-\bar\alpha_t)I
    \right),
    \quad
    \text{where} \quad
    \alpha_t := 1-\beta_t,
    \qquad
    \bar\alpha_t := \prod_{s=1}^{t}\alpha_s.
    \label{eq:ddpm_forward_marginal}
\end{equation}
Equivalently, \(x_t\) can be written through the reparameterization
\begin{equation}
    x_t
    =
    \sqrt{\bar\alpha_t}\,x_0
    +
    \sqrt{1-\bar\alpha_t}\,\varepsilon,
    \qquad
    \varepsilon \sim \mathcal{N}(0, I).
    \label{eq:ddpm_forward_reparam}
\end{equation}
This makes explicit that the state at timestep \(t\) is a linear combination of the original sample and Gaussian noise. As \(t\) increases, the signal coefficient \(\sqrt{\bar\alpha_t}\) decreases, so the forward process progressively destroys structure and drives \(x_t\) toward an isotropic Gaussian distribution.

%%%%%%%%%%%%%%%%%%%%%%%%%%%%%%%%%%%%%%%%%%%%%%%%%%%%%%%%%%%%%%%%%%%%%%%%%%%%%%%
\subparagraph{Reverse Denoising Process.}
%%%%%%%%%%%%%%%%%%%%%%%%%%%%%%%%%%%%%%%%%%%%%%%%%%%%%%%%%%%%%%%%%%%%%%%%%%%%%%%

The generative task is to invert the forward corruption process. To achieve this, \citet{NEURIPS2020_4c5bcfec} learn a reverse Markov chain whose transitions are Gaussian:
\begin{equation}
    p_\theta(x_{t-1}\mid x_t)
    =
    \mathcal{N}\!\left(
        \mu_\theta(x_t,t),\;
        \Sigma_\theta(x_t,t)
    \right).
    \label{eq:ddpm_reverse}
\end{equation}
Here, \(\theta\) denotes the parameters of a single neural network shared across timesteps, and the timestep \(t\) is provided to the network through a timestep embedding.

Figure~\ref{fig:variational_diffusion_model} gives an overall picture of the diffusion process, showing that the forward and reverse chains pass through the same sequence of latent states in opposite directions.
\begin{figure}[ht]
    \centering
    \includegraphics[width=0.45\linewidth]{chapters/background/media/variational_diffusion_model.png}
    \caption{\citet{variational_diffusion_model} illustrate a variational diffusion model in which the lower path represents the forward process that gradually corrupts \(x_0\) into \(x_T\), while the upper path represents the reverse process that denoises step by step. 
    The intermediate states \(x_1,\dots,x_{T-1}\) are latent variables of the same dimensionality as the data \citep{tutorial_on_diffusion_models}.}
    \label{fig:variational_diffusion_model}
\end{figure}

Because the forward process is Gaussian and fixed, it induces an exact conditional distribution over the previous state when both the current noisy state \(x_t\) and the original clean sample \(x_0\) are known:
\begin{equation}
    q(x_{t-1}\mid x_t,x_0)
    =
    \mathcal{N}\!\left(
        \tilde{\mu}_t(x_t,x_0),\;
        \tilde{\beta}_t I
    \right).
    \label{eq:true_reverse_posterior}
\end{equation}
Its mean and variance are
\begin{align}
    \tilde{\mu}_t(x_t,x_0)
    &=
    \frac{\sqrt{\bar\alpha_{t-1}}\beta_t}{1-\bar\alpha_t}\,x_0
    +
    \frac{\sqrt{\alpha_t}(1-\bar\alpha_{t-1})}{1-\bar\alpha_t}\,x_t,
    \label{eq:true_reverse_mean}
    \\
    \tilde{\beta}_t
    &=
    \frac{1-\bar\alpha_{t-1}}{1-\bar\alpha_t}\,\beta_t.
    \label{eq:true_reverse_var}
\end{align}

It is important to distinguish the reverse conditional in Equation~\ref{eq:true_reverse_posterior} from the forward transition in Equation~\ref{eq:ddpm_forward_transition}.
Since the forward process is Markovian,
\begin{equation}
    q(x_t\mid x_{t-1},x_0)=q(x_t\mid x_{t-1}),
    \label{eq:forward_markov}
\end{equation}
so the next noisy state depends only on the previous noisy state. 
In contrast, the exact reverse conditional in Equation~\ref{eq:true_reverse_posterior} depends on both \(x_t\) and \(x_0\). During training, \(x_0\) is known, so this conditional can be computed exactly. At sampling time, however, \(x_0\) is unknown, so the model learns the practical approximation \(p_\theta(x_{t-1}\mid x_t)\) in Equation~\ref{eq:ddpm_reverse}.

Figure~\ref{fig:transition_block} highlights this local asymmetry. The forward transition is fixed and Gaussian, while the reverse transition must be learned.

\begin{figure}[ht]
    \centering
    \includegraphics[width=0.60\linewidth]{chapters/background/media/transition_block.png}
    \caption{Local transition block in a diffusion model. 
    The blue circles denote neighboring latent states in the diffusion chain: 
    \(x_{t-1}\) is a less noisy state, \(x_t\) is the current state, and \(x_{t+1}\) is a more noisy state. 
    The forward process moves from left to right, gradually adding noise, while the reverse process moves from right to left, progressively denoising. 
    The forward transition \(q(x_t\mid x_{t-1})\) is the fixed Gaussian noising step that maps a cleaner state to a noisier one. 
    The learned reverse transition \(p_\theta(x_{t-1}\mid x_t)\) maps a noisy state to a less noisy one and is used at sampling time. 
    The quantity \(p(x_{t-1}\mid x_t)\) denotes the ideal but generally intractable reverse transition, while the exact reverse conditional \(q(x_{t-1}\mid x_t,x_0)\) becomes tractable during training because the original clean sample \(x_0\) is known. 
    Thus, the forward process uses the fixed Gaussian transition \(q(x_t\mid x_{t-1})\), while the reverse process learns \(p_\theta(x_{t-1}\mid x_t)\). 
    The exact reverse conditional \(q(x_{t-1}\mid x_t,x_0)\) is available during training because \(x_0\) is known, but not during sampling \citep{tutorial_on_diffusion_models}.}
    \label{fig:transition_block}
\end{figure}

%%%%%%%%%%%%%%%%%%%%%%%%%%%%%%%%%%%%%%%%%%%%%%%%%%%%%%%%%%%%%%%%%%%%%%%%%%%%%%%
\subparagraph{Training Objective.}
%%%%%%%%%%%%%%%%%%%%%%%%%%%%%%%%%%%%%%%%%%%%%%%%%%%%%%%%%%%%%%%%%%%%%%%%%%%%%%%

Training seeks to maximize the likelihood of the data under the learned reverse process. 
Since the latent variables \(x_{1:T}\) cannot be marginalized exactly, training instead optimizes a standard variational lower bound:
\begin{equation}
    \mathbb{E}_q[-\log p_\theta(x_0)]
    \;\le\;
    \mathcal{L}(\theta)
    =
    \mathbb{E}_q\!\left[
        -\log\frac{p_\theta(x_{0:T})}{q(x_{1:T}\mid x_0)}
    \right].
    \label{eq:elbo_ddpm_raw}
\end{equation}
This can be rewritten as the familiar decomposition
\begin{equation}
    \mathcal{L}(\theta)
    =
    \mathbb{E}_q\Big[
        \underbrace{
            D_{\mathrm{KL}}\bigl(q(x_T\mid x_0)\,\|\,p(x_T)\bigr)
        }_{L_T}
        +
        \sum_{t=2}^{T}
        \underbrace{
            D_{\mathrm{KL}}\bigl(
                q(x_{t-1}\mid x_t,x_0)
                \;\|\;
                p_\theta(x_{t-1}\mid x_t)
            \bigr)
        }_{L_{t-1}}
        -
        \underbrace{\log p_\theta(x_0\mid x_1)}_{L_0}
    \Big].
    \label{eq:elbo_ddpm_kl}
\end{equation}
The term \(L_T\) matches the terminal forward distribution to the prior. The intermediate terms \(L_{t-1}\) compare the learned reverse transition in Equation~\ref{eq:ddpm_reverse} with the exact conditional in Equation~\ref{eq:true_reverse_posterior}. The final term \(L_0\) plays the role of a reconstruction term for \(x_0\).

%%%%%%%%%%%%%%%%%%%%%%%%%%%%%%%%%%%%%%%%%%%%%%%%%%%%%%%%%%%%%%%%%%%%%%%%%%%%%%%
\subparagraph{Simplified Noise-Prediction Objective.}
%%%%%%%%%%%%%%%%%%%%%%%%%%%%%%%%%%%%%%%%%%%%%%%%%%%%%%%%%%%%%%%%%%%%%%%%%%%%%%%

The variational bound in Equation~\ref{eq:elbo_ddpm_kl} decomposes as
\(
    L_{\mathrm{VLB}} = L_T + \sum_{t=2}^T L_{t-1} + L_0.
\)
The terms \(L_T\) and \(L_0\) do not depend on how the learned reverse mean is parameterized, so the key modeling choice appears in the intermediate terms \(L_{t-1}\).

Since both \(q(x_{t-1}\mid x_t,x_0)\) and \(p_\theta(x_{t-1}\mid x_t)\) are Gaussian, and the reverse-process variance is fixed, each intermediate term reduces to matching the learned mean \(\mu_\theta(x_t,t)\) to the true posterior mean \(\tilde{\mu}_t(x_t,x_0)\):
\begin{equation}
    L_{t-1}
    =
    \mathbb{E}_q\!\left[
        \frac{1}{2\sigma_t^2}
        \bigl\|
            \tilde{\mu}_t(x_t,x_0)
            -
            \mu_\theta(x_t,t)
        \bigr\|^2
    \right]
    + C_t,
    \label{eq:L_t_mse_form}
\end{equation}
where \(C_t\) does not depend on \(\theta\).

Using the forward reparameterization in Equation~\ref{eq:ddpm_forward_reparam}, the true posterior mean in Equation~\ref{eq:true_reverse_mean} can be rewritten as
\begin{equation}
    \tilde{\mu}_t(x_t,x_0)
    =
    \frac{1}{\sqrt{\alpha_t}}
    \left(
        x_t
        -
        \frac{1-\alpha_t}{\sqrt{1-\bar{\alpha}_t}}\,
        \varepsilon
    \right).
    \label{eq:true_reverse_mean_eps}
\end{equation}
This suggests a convenient parameterization of the learned reverse mean. Instead of predicting \(\mu_\theta(x_t,t)\) directly, the network predicts the noise term \(\varepsilon_\theta(x_t,t)\), and the model mean is then defined by
\begin{equation}
    \mu_\theta(x_t,t)
    =
    \frac{1}{\sqrt{\alpha_t}}
    \left(
        x_t
        -
        \frac{1-\alpha_t}{\sqrt{1-\bar{\alpha}_t}}\,
        \varepsilon_\theta(x_t,t)
    \right).
    \label{eq:reverse_mean_eps_param}
\end{equation}
Thus, the network is trained to predict noise, while this prediction implicitly determines the learned reverse mean.

Substituting Equations~\ref{eq:true_reverse_mean_eps} and \ref{eq:reverse_mean_eps_param} into Equation~\ref{eq:L_t_mse_form} gives
\begin{equation}
    L_{t-1}
    =
    \mathbb{E}_{x_0,\varepsilon}
    \left[
        \frac{(1-\alpha_t)^2}{2\sigma_t^2\,\alpha_t(1-\bar{\alpha}_t)}
        \,
        \bigl\|
            \varepsilon
            -
            \varepsilon_\theta(x_t,t)
        \bigr\|^2
    \right]
    + C_t,
    \label{eq:weighted_noise_loss}
\end{equation}
so each intermediate KL term becomes a weighted mean-squared error between the true noise \(\varepsilon\) and the noise predicted by the network.

Dropping the timestep-dependent weight yields the simplified objective used in practice:
\begin{equation}
    L_{\mathrm{simple}}(\theta)
    =
    \mathbb{E}_{t,x_0,\varepsilon}
    \bigl[
        \|
            \varepsilon
            -
            \varepsilon_\theta(x_t,t)
        \|^2
    \bigr],
    \qquad
    x_t
    =
    \sqrt{\bar{\alpha}_t}\,x_0
    +
    \sqrt{1-\bar{\alpha}_t}\,\varepsilon.
    \label{eq:ddpm_simple_loss}
\end{equation}
This is a reweighted form of the ELBO, but is simpler to optimize and is commonly used because it gives stable training and high-quality samples.

%%%%%%%%%%%%%%%%%%%%%%%%%%%%%%%%%%%%%%%%%%%%%%%%%%%%%%%%%%%%%%%%%%%%%%%%%%%%%%%
\subparagraph{Training and Sampling Algorithms.}
%%%%%%%%%%%%%%%%%%%%%%%%%%%%%%%%%%%%%%%%%%%%%%%%%%%%%%%%%%%%%%%%%%%%%%%%%%%%%%%

Training and sampling use the same network in different ways. 
During training, the model does not execute the full reverse chain. 
Instead, it draws a clean sample \(x_0\), samples a timestep \(t\) uniformly from \(\{1,\dots,T\}\), samples noise \(\varepsilon\sim\mathcal{N}(0,I)\), and constructs \(x_t\) directly using Equation~\ref{eq:ddpm_forward_reparam}. 
The network then predicts \(\varepsilon_\theta(x_t,t)\), and a gradient step is taken on the single-timestep loss in Equation~\ref{eq:ddpm_simple_loss}. 
Thus, the gradient is computed from the loss at the sampled noise level \(t\); training does not backpropagate through the full reverse trajectory.

This still trains one model to handle all timesteps because the same network is shared across the entire noise schedule and receives \(t\) as an input. 
By sampling different values of \(t\) during training, the model learns to denoise across all noise levels.

The full sequential reverse chain appears only during sampling. 
Starting from \(x_T\sim\mathcal{N}(0,I)\), the model repeatedly applies the reverse update
\begin{equation}
    x_{t-1}
    =
    \frac{1}{\sqrt{\alpha_t}}
    \left(
        x_t
        -
        \frac{\beta_t}{\sqrt{1-\bar\alpha_t}}\,
        \varepsilon_\theta(x_t,t)
    \right)
    +
    \sigma_t z,
    \qquad
    z\sim\mathcal{N}(0,I),
\end{equation}
for \(t=T,T-1,\dots,1\), where \(\sigma_t^2\) is fixed by the chosen reverse-process parameterization. 
This iterative reverse process gradually transforms noise into a structured sample \(x_0\).
% %%%%%%%%%%%%%%%%%%%%%%%%%%%%%%%%%%%%%%%%%%%%%%%%%%%%%%%%%%%%%%%%%%%%%%%%%%%%%%%
% \subparagraph{Architecture.}
% %%%%%%%%%%%%%%%%%%%%%%%%%%%%%%%%%%%%%%%%%%%%%%%%%%%%%%%%%%%%%%%%%%%%%%%%%%%%%%%

% The reverse-process predictor is typically implemented with a U-Net architecture with residual blocks, skip connections, and multi-scale feature hierarchies \citep{NEURIPS2020_4c5bcfec}. The timestep \(t\) is encoded by a sinusoidal positional embedding and injected into intermediate layers so that the network can adapt its prediction to the current noise level. The network outputs \(\varepsilon_\theta(x_t,t)\), which is converted to the reverse mean through Equation~\ref{eq:reverse_mean_eps_param}. In standard DDPM implementations, the reverse variance is fixed and not predicted by the network.

\subsection{Score-Based Formulation}

Song et al.~\citep{song2021scorebased} showed that DDPMs can be understood in continuous time
through stochastic differential equations. In this view, discrete-time diffusion models and
score-based generative models are two parameterizations of the same underlying denoising process.

\begin{figure}[H]
    \centering
    \includegraphics[width=0.45\textwidth]{chapters/background/media/score.jpg}
    \caption{\textbf{Forward and reverse SDEs.} In continuous time, the forward process corrupts
    data with drift and diffusion terms, while the reverse process uses the score function to
    steer noisy samples back toward the data manifold \citep{song2021scorebased}.}
    \label{fig:sde_score_diagram}
\end{figure}

\paragraph{From the Markov Chain to the Continuous Limit.}
The discrete forward chain can be interpreted as a discretized version of a continuous corruption process. In continuous time, the variance-preserving formulation is described by
\begin{equation}
    \mathrm{d}x
    =
    -\tfrac{1}{2}\beta(t)\,x\,\mathrm{d}t
    +
    \sqrt{\beta(t)}\,\mathrm{d}W(t),
\end{equation}
where $W(t)$ is a Wiener process. More generally, the forward corruption process is written as
\begin{equation}
    \mathrm{d}x = f(x,t)\,\mathrm{d}t + g(t)\,\mathrm{d}w,
\end{equation}
with drift \(f(x,t)\) and diffusion strength \(g(t)\).

\paragraph{Reverse SDE.}
The corresponding reverse-time SDE has the form
\begin{equation}
    \mathrm{d}x
    =
    \left[
        f(x,t) - g(t)^2 \nabla_x \log p_t(x)
    \right]\mathrm{d}t
    +
    g(t)\,\mathrm{d}\bar w,
\end{equation}
where \(\nabla_x \log p_t(x)\) is the \emph{score function}, namely the gradient of the
log-density at noise level \(t\). Intuitively, the score points in the direction that makes
the current sample more likely under the time-dependent data distribution. The reverse SDE is
the continuous-time counterpart of the DDPM reverse chain where reverse denoising is driven by a vector field
that points back toward the data manifold.

In the DDPM parameterization, the model outputs \(\varepsilon_\theta(x_t,t)\), an estimate
of the Gaussian noise used to corrupt \(x_0\) into \(x_t\). Because \(x_t\) is a linear
combination of signal and noise, this noise estimate can be converted directly into a score
estimate:
\begin{equation}
    s_\theta(x_t, t) = -\frac{\varepsilon_\theta(x_t, t)}{\sqrt{1-\bar\alpha_t}}.
    \label{eq:score_noise}
\end{equation}
DDPMs can therefore be viewed as score-based generative models written in a different
parameterization. 
% This connection is especially relevant for Phase~1, in Section~\ref{sec:phase_1}, where composition is performed by combining score estimates from pretrained score-based diffusion models during inference.


\subsection{Conditional Generation and Classifier-Free Guidance}

This section distinguishes between \textbf{unconditional} and \textbf{conditional} diffusion models by examining their objectives, formulations, and applications.

While the above formulation defines the foundation of \emph{unconditional diffusion models}, practical applications often demand controlled or semantically guided generation such as synthesizing specific classes, attributes, or text-aligned imagery.
This introduces the need for \emph{conditional diffusion models}, where additional information \(y\) (e.g., labels, text embeddings, or medical attributes) guides the reverse process.
Recent extensions, such as \emph{classifier} \citep{NEURIPS2021_49ad23d1} and \emph{classifier-free guidance} \citep{ho2022classifier}, further refine this conditioning by explicitly or implicitly modulating the score function to balance fidelity and diversity during sampling.

\paragraph{Unconditional Diffusion Models.}
An \textit{unconditional diffusion model} learns to generate samples from the marginal data distribution \(p(x)\) without using auxiliary information or labels. During training, the model observes only samples \(x_0\) drawn from \(p_{\text{data}}(x)\) and is optimized to predict the noise added at each diffusion step:
\begin{equation}
\mathcal{L}_{\text{uncond}} = \mathbb{E}_{x_0, t, \varepsilon}\!\left[\|\varepsilon - \varepsilon_\theta(x_t, t)\|^2\right].
\end{equation}
The training signal is generated from the data itself rather than supplied by external annotation, so the objective is self-supervised. At sampling time, the model starts from Gaussian noise \(x_T \sim \mathcal{N}(0, I)\) and denoises it step by step to produce a realistic sample. In this sense, unconditional diffusion models learn \(p(x)\) directly, without any conditioning variable.

\paragraph{Conditional Diffusion Models.}
A \textit{conditional diffusion model} extends the unconditional setting by modeling
\(p(x \mid y)\), where \(y\) provides additional information such as a class label, text
prompt, or spatial map. The forward noising process is unchanged. Conditioning enters only
through the reverse model, which is now written as
\[
p_\theta(x_{t-1}\mid x_t,y).
\]
In the DDPM parameterization, this is implemented by conditioning the noise predictor
\(\varepsilon_\theta(x_t,t,y)\). Training therefore uses the objective
\begin{equation}
\mathcal{L}_{\text{cond}}
=
\mathbb{E}_{x_0,t,\varepsilon,y}
\left[
\|\varepsilon-\varepsilon_\theta(x_t,t,y)\|^2
\right].
\end{equation}
By Equation~\ref{eq:score_noise}, this is equivalent to learning the conditional score of
the noisy data distribution. At sampling time, the condition \(y\) steers the denoising
trajectory toward samples consistent with the desired class, prompt, or attribute. In this
way, conditional diffusion turns unconditional generation into controllable generation.

\paragraph{Classifier Guidance.}
Classifier guidance \citep{NEURIPS2021_49ad23d1} provides a post-hoc way to steer an
unconditional diffusion model toward a target condition \(y\). The key observation is that
a classifier gives access to \(p_\phi(y\mid x_t)\), and therefore to the gradient
\[
\nabla_{x_t}\log p_\phi(y\mid x_t).
\]
At timestep \(t\), the goal is to obtain the conditional score of the noisy-state distribution \(p_t(x_t\mid y)\). Using the product rule for the joint distribution,
\[
p_t(x_t,y)=p_\phi(y\mid x_t)\,p_t(x_t)=p_t(x_t\mid y)\,p(y).
\]
Rearranging the terms yields
\[
p_t(x_t\mid y)=\frac{p_\phi(y\mid x_t)\,p_t(x_t)}{p(y)}.
\]
Therefore,
\[
\log p_t(x_t\mid y)
=
\log p_\phi(y\mid x_t)+\log p_t(x_t)-\log p(y),
\]
and differentiating with respect to \(x_t\) yields
\[
\nabla_{x_t}\log p_t(x_t\mid y)
=
\nabla_{x_t}\log p_t(x_t)
+
\nabla_{x_t}\log p_\phi(y\mid x_t),
\]
since \(p(y)\) does not depend on \(x_t\).The conditional score can therefore be written
as the sum of the unconditional score and a classifier-derived conditioning term. In
practice, \citet{NEURIPS2021_49ad23d1} scale this conditioning term by a guidance weight
\(\gamma \ge 0\):
\begin{equation}
\tilde{s}(x_t,t,y)
=
s_\theta(x_t,t)
+
\gamma\,\nabla_{x_t}\log p_\phi(y\mid x_t).
\label{eq:classifier_guidance_scaled}
\end{equation}
Larger values of \(\gamma\) strengthen adherence to the target condition, often at the cost
of reduced sample diversity.

\paragraph{Classifier-Free Guidance.}
Classifier-free guidance (CFG) \citep{ho2022classifier} achieves a similar effect without
training a separate classifier. A single diffusion model is trained with condition dropout
so that it learns both the conditional and unconditional predictors
\[
\varepsilon_\theta(x_t,t,y)
\qquad\text{and}\qquad
\varepsilon_\theta(x_t,t,\emptyset),
\]
where \(\emptyset\) denotes the null condition. At sampling time, these two predictions are
combined as
\begin{equation}
\tilde{\varepsilon}_\theta(x_t,t,y)
=
(1+w)\,\varepsilon_\theta(x_t,t,y)
-
w\,\varepsilon_\theta(x_t,t,\emptyset),
\label{eq:guided_eps_cfg}
\end{equation}
with guidance weight \(w \ge 0\). When \(w=0\), this reduces to ordinary conditional
generation. Larger values of \(w\) strengthen the conditional signal while reducing
diversity. This formulation has become the standard conditioning mechanism in modern
diffusion models. In the proposed study, it plays two concrete roles. 
This is relevant throughout the proposal because classifier-free guidance provides the basic conditioning mechanism used for standard prompting and for composition. 
In Phase~1, it underlies both the single-prompt baseline and the compositional operator of \citet{liu2022compositional}. 
In Phases~2 and~3, the same conditional prediction framework is used in the controlled synthetic and clinical composition settings.

\paragraph{Stable Diffusion.}
Stable Diffusion \citep{rombach2022high} is a \emph{latent diffusion model}, meaning that diffusion is performed in a compressed latent space rather than directly in pixel space. 
An image \(x\) is first mapped by a variational autoencoder into a latent representation \(z\), and the forward and reverse diffusion processes are learned over this lower-dimensional latent variable. T
his substantially reduces the computational cost of training and sampling while retaining the ability to decode high-resolution images. 
For text-to-image generation, a text encoder produces prompt embeddings that are injected into the U-Net denoiser through cross-attention \citep{vaswani2017attention}, allowing the denoising trajectory to remain aligned with the prompt semantics. 
During training, the model learns to predict the noise added to latent variables under text conditioning. 
During sampling, it starts from Gaussian latent noise, iteratively denoises using the text-conditioned predictor with classifier-free guidance, and then decodes the final latent through the VAE decoder to obtain an image. 
\begin{figure}[H]
    \centering
    \includegraphics[width=0.60\textwidth]{chapters/background/media/stable diffusion.png}
    \caption{Stable Diffusion as a latent text-to-image diffusion model. A VAE maps images to a compressed latent space, a text encoder provides prompt embeddings, and a U-Net denoises latent noise under text conditioning before the VAE decoder reconstructs the final image. Operating in latent space reduces computational cost while preserving high-quality generation.}
    \label{fig:stable_diffusion_architecture}
\end{figure}

\subsection{Vision-Language Models and Prompt Inversion}

\paragraph{CLIP.}
CLIP \citep{radford2019language} learns a shared embedding space for images and text using a contrastive objective over paired image-caption data. 
Its central contribution is to align visual and textual representations so that matching image-text pairs have high similarity while mismatched pairs are pushed apart. 
This shared space is useful for text-to-image systems because it provides a generic notion of semantic correspondence between prompts and images. 
In the context of the proposed study, CLIP is relevant as a representation for prompt-image matching and as a building block for inversion methods that map generated images back toward text-aligned embeddings.

\paragraph{BLIP.} BLIP \citep{li2022blip} is a vision-language model for both image understanding and language generation. 
It combines image and text encoders with multimodal pre-training to support tasks such as captioning, retrieval, and visual question answering. 
Unlike contrastive models such as CLIP, which mainly return cross-modal similarity scores, BLIP can generate text conditioned on an image. This makes it useful when the evaluation pipeline requires a textual description of generated outputs rather than only an image-text matching score.

BLIP-VQA denotes the use of BLIP-style vision-language models for visual question answering, where the model is conditioned on both an image and a natural-language question and returns a textual answer. 
In this setting, the model can be asked whether particular entities, attributes, or relations are present in a generated image, making it useful for targeted evaluation of compositional correctness. 
For the proposed study, BLIP-VQA is relevant as a possible tool for querying whether a generated sample actually contains the intended concepts, rather than relying only on global similarity measures or free-form captions.

\paragraph{Prompt Inversion.}
In the CLIP \citep{radford2019language} text encoder used by Stable Diffusion \citep{rombach2022high, podell2023sdxl}, the End-of-Sentence (EOS) token has an important role beyond simply marking the end of a prompt. 
Because CLIP uses a causal attention mask, each token can attend only to the tokens that come before it. 
As the final token in the sequence, the EOS position is therefore the only one whose hidden state has access to the entire prompt. 
When CLIP forms a single text representation for image-text alignment, it uses this EOS hidden state rather than averaging all token embeddings. 
In this sense, the EOS representation acts as a contextualized summary of the whole prompt.

Stable Diffusion Image-to-Prompt Conversion (SD-IPC) \citep{the_clip_model_is_secretly_an_image_to_prompt_converter} builds on this property to convert a reference image into a Stable Diffusion text-conditioning signal with little additional training. 
The method relies on two observations: first, that Stable Diffusion generation is strongly influenced by this EOS-based text representation; and second, that CLIP training aligns the projected EOS text embedding with the corresponding visual embedding. 
Using this relationship, the authors derive a closed-form pseudo-inverse of the CLIP text projection matrix to map a CLIP image embedding into the prompt-side conditioning space. 
This converted embedding is then repeated to form a pseudo-prompt sequence. 
In the context of this study, SD-IPC is useful because it provides a concrete way to test whether a target image can be recovered from Stable Diffusion's prompt-conditioning space.